% !TEX root = ../main.tex

\chapter{Introducción al Álgebra Lineal}
\label{cap:introduccion}

Bienvenido al primer capítulo. Fíjate en la primera línea de este archivo: ese comentario le dice a VS Code que el documento real que debe compilar está un nivel de carpetas más arriba y se llama main.tex.

\section{Espacios Vectoriales}
\label{sec:espacios_vectoriales}

Comenzaremos definiendo uno de los conceptos fundamentales. Usaremos los entornos y comandos que definimos en nuestro archivo de estilos.

\begin{definicion}[Espacio Vectorial]
    \label{def:espacio_vectorial}
    Un espacio vectorial sobre un cuerpo (como $\R$ o $\C$) es un conjunto no vacío dotado de dos operaciones...
\end{definicion}

Fíjate cómo usamos nuestras macros $\R$ y $\C$ directamente en el texto. Ahora, enunciemos un resultado clásico utilizando la norma y el valor absoluto que configuramos con mathtools.

\begin{teorema}[Desigualdad de Cauchy-Schwarz]
    \label{teo:cauchy}
    Para todo par de vectores $u, v$ en un espacio prehilbertiano, se cumple que:
    $$ \abs{\langle u, v \rangle} \leq \norm{u} \norm{v} $$
\end{teorema}

\begin{proof}
    La demostración de este resultado se basa en construir un polinomio cuadrático y analizar su discriminante...
\end{proof}

Como vimos en el Teorema \ref{teo:cauchy}, la relación entre la norma y el producto interno es la base de la geometría en estos espacios.